Domain randomisation trains a policy on many simulated variants of a task so that it tolerates the mismatch to a target system. It was popularised for visual transfer \cite{tobin2017domain} and for dynamics transfer in robot control \cite{peng2018sim}, and it was a key ingredient of the dexterous-manipulation system of \cite{openai2019solving}, whose Automatic Domain Randomization widens ranges as the policy succeeds. Dynamics randomisation is not always needed, however: Xie et al.\ report direct transfer of a quadruped locomotion policy without it \cite{xie2021dynamics}. A practical question is rarely answered explicitly: how wide should the randomisation ranges be? Too narrow and the policy breaks under shift; too wide and, it is often argued that the policy becomes conservative or fails to learn, so nominal performance suffers (we did not find a source that measures this cost directly). Theory on DR formalises the dependence of transfer on the randomisation distribution \cite{chen2021understanding}, but practitioners still set ranges by hand or adapt them with real data \cite{chebotar2019closing,tiboni2022dropo}.

We ask a narrow, falsifiable version of the question in a setting cheap enough to run many seeds: is there a range width beyond which nominal return drops faster than out-of-nominal robustness improves? We also test whether a simple automatic curriculum that widens the range on success, an instance of the curricula surveyed in \cite{portelas2020automatic}, avoids having to choose a width.

\textbf{Contributions.} (i) A controlled sweep of DR width on CartPole with a CEM-trained policy, evaluated on nominal, in-range and out-of-range dynamics (5 seeds); (ii) a comparison of no DR, narrow DR, wide DR and a success-gated curriculum, with ablations of the gate; (iii) an honest negative finding: in this setting we did not observe a nominal-performance cost of wide randomisation, and the gate was not useful. We make no claim beyond this small setting.
